\section{The \texorpdfstring{$\STS(15)$}{STS(15)} defect spectrum}\label{sec:app-sts}

This appendix records the complete realized defect spectrum on $15$ points, the
general squarefree trichotomy, and the ternary augmentation splitting for triple
designs. The first two follow from the main results
and classical $p$-rank theorems; the third is proved here and is used by the
ternary case of Appendix~\ref{sec:app-er}.

\subsection{The spectrum and primary purity}

\begin{theorem}[the $\STS(15)$ spectrum]\label{thm:sts15}
For a Steiner triple system on $15$ points, $\rank_{\FF_3}(A)=14$ and
$\rank_{\FF_3}(\Cx)=15$; hence $\Delt=(\ZZ/2)^{a}$ with $a=15-\rank_{\FF_2}(A)$.
The realized binary ranks are $11,12,13,14,15$~\cite{AssmusKey1992}, so the
complete realized spectrum is
\[
\Delt\in\{0,\ \ZZ/2,\ (\ZZ/2)^{2},\ (\ZZ/2)^{3},\ (\ZZ/2)^{4}\},
\]
with full incidence forms $1^{14}3,\ 1^{14}6,\ 1^{13}2\cdot6,\ 1^{12}2^{2}6,\
1^{11}2^{3}6$ respectively. More generally, using that a nontrivial $\STS(v)$ has
full $p$-rank for every prime $p\ge5$~\cite{DoyenHubautVandensavel} together with
the no-simultaneous-deficiency theorem of Shi--Xu--Krotov~\cite{ShiXuKrotov},
every nontrivial $\STS$ defect is a pure $2$-group, a pure $3$-group, or zero.
\end{theorem}

The $\STS(15)$ ternary statement is structural (a nonconstant ternary dual
vector would partition the points into three $\STS(5)$'s, impossible); the five
binary ranks are imported from Assmus--Key~\cite{AssmusKey1992} and
Tonchev--Weishaar~\cite{TonchevWeishaar}, and the code/affine-lattice
classification of the $80$ systems is Sol\'e's~\cite{Sole1997}. The general
trichotomy translates the cited $p$-rank results into the selected quotient.

\subsection{Removing the ternary augmentation factor}

\begin{proposition}[ternary augmentation splitting]\label{prop:split}
For every nontrivial $2$-$(v,3,\lambda)$ design with $v>3$, the all-ones block
row is not in the $\FF_3$-row span of the incidence matrix $A\colon\ZZ^{\mathcal
B}\to\ZZ^{P}$; consequently the augmentation
$\varepsilon\colon\coker(A)\to\ZZ/3$ splits, so
$\coker(A)\cong\Delt(D)\oplus\ZZ/3$.
\end{proposition}

\begin{proof}
Suppose $a\in\FF_3^{P}$ satisfies $\sum_{x\in B}a_x=1$ for every block $B$, and
partition the points by fibre $P_i=\{x:a_x=i\}$, $n_i=|P_i|$ ($i\in\ZZ/3$). A
triple sums to $1$ in $\FF_3$ only if its values form the multiset $(i,i,i+1)$;
call such a block \emph{type $i$}. Each pair inside $P_i$ lies in $\lambda$
blocks and forces the third point into $P_{i+1}$, giving $\lambda\binom{n_i}{2}$
blocks of type $i$; counting the same blocks by their two $P_i$--$P_{i+1}$
cross-pairs gives $2\cdot(\#\,\text{type }i)=\lambda\,n_i n_{i+1}$. Hence
$\lambda\,n_i(n_i-1)=\lambda\,n_i n_{i+1}$, so $n_{i+1}=n_i-1$ whenever $n_i>0$.
Applied cyclically over $i\in\ZZ/3$ this forces three strict decreases
$n_0>n_1>n_2>n_0$, impossible unless some fibre is empty; and with an empty fibre
the surviving instances of $n_{i+1}=n_i-1$ bound $n_0+n_1+n_2\le3$ (the maximum
is the rotation $(0,2,1)$), contradicting $v>3$.

Thus no such $a$ exists, so the
all-ones row lies outside the $\FF_3$-row span of $A$.

Row-space duality then
supplies $y\in\FF_3^{\mathcal B}$ with $Ay=0$ in $\FF_3^{P}$ and
$\mathbf 1_{\mathcal B}^{T}y=1$. Lift $y$ to $\tilde y\in\ZZ^{\mathcal B}$; then
$A\tilde y\equiv0\pmod 3$, so $z:=A\tilde y/3\in\ZZ^{P}$ is integral, its class
$[z]\in\coker(A)$ satisfies $3[z]=[A\tilde y]=0$, and, using
$\mathbf 1_P^{T}A=3\,\mathbf 1_{\mathcal B}^{T}$ (each block has three points),
\[
\varepsilon([z])=\mathbf 1_P^{T}z\bmod 3
=\tfrac{1}{3}\mathbf 1_P^{T}A\tilde y\bmod 3
=\mathbf 1_{\mathcal B}^{T}\tilde y\bmod 3=1 .
\]
Hence $\ZZ/3\to\coker(A)$, $1\mapsto[z]$, is a section of $\varepsilon$, and
$\coker(A)\cong\Delt(D)\oplus\ZZ/3$.
\end{proof}
